\section{Conclusion}
\label{sec:conclusion}

This paper has developed a five-type taxonomy of state redistricting regimes---permissive, criteria-based, commission, hybrid, and restrictive---and demonstrated that a single METIS recursive bisection algorithm, parameterized through per-state YAML configuration files, can satisfy the dominant legal criteria across all five types.

The key findings are threefold. First, the mapping from legal criteria to algorithmic parameters is comprehensive: compactness requirements map to edge-weight scaling, county-preservation requirements map to adjacency weight overrides, community-of-interest requirements map to custom edge-weight matrices, and partisan-neutrality requirements are satisfied structurally through the algorithm's architectural inability to receive partisan data. No state redistricting criterion examined in this paper requires an algorithmic capability that is not already present in the METIS framework.

Second, the algorithm's structural partisan blindness satisfies partisan-neutrality requirements more completely than any procedural approach. States like Florida (Amendment 6) and Michigan (Proposal 2) prohibit partisan consideration in redistricting; the algorithmic approach satisfies this prohibition categorically and verifiably, not merely on the basis of mapmaker promises. This categorical satisfaction is a stronger constitutional claim and is independently verifiable through code audit.

Third, the five-type taxonomy reveals that approximately one-third of states (Types I and II) can be served by straightforward parameterization of the default algorithm, while the remaining two-thirds (Types III--V) require additional configuration for compactness, community of interest, and partisan neutrality. The additional configuration is achievable through existing parameter channels without any modification to the core algorithm.

The practical implication is that a national algorithmic redistricting system need not be designed as 50 separate algorithms. A single architecture with a 50-entry YAML configuration table can serve all states, adapting to legal heterogeneity through parameter variation. This represents a significant simplification of the implementation challenge and preserves the reproducibility and transparency advantages of the algorithmic approach.

Future work should address three open questions. First, the community-of-interest mapping is underdeveloped: the current approach uses census-designated places as a proxy, but actual communities of interest are more granular and context-specific. A richer operationalization---perhaps using commuting-zone data, school district boundaries, or demographic similarity metrics---would improve the algorithm's performance on this criterion. Second, the competitive-elections criterion (Arizona) requires further empirical analysis to establish whether METIS maps systematically produce more competitive districts than enacted maps across multiple census years and redistricting cycles. Third, the legal question of whether algorithmic maps satisfy the ``intent'' prong of Florida's Amendment 6 has not been adjudicated; a definitive answer will require either legislative adoption or test litigation.

The broader conclusion is that the legal heterogeneity of state redistricting criteria, while real and significant, does not present an obstacle to national algorithmic adoption. The algorithm is more adaptable than it appears when viewed from a single state's legal perspective.
